\chapter{Cross Sections, Phase Space, and Unitarity}
\label{ch:cross-sections-unitarity}

\section{Transition Probabilities}

Let \(\mathcal H_{\rm as}\) be the asymptotic Hilbert space in a sector where
the incoming and outgoing Haag--Ruelle ranges coincide.  The scattering
operator
\[
  S:\mathcal H_{\rm as}\to\mathcal H_{\rm as}
\]
is then unitary.  Write
\[
  S=\mathbf 1+\ii T .
\]
For a normalized incoming vector \(\Psi_{\rm in}\in\mathcal H_{\rm as}\) and
an outgoing class of states \(\mathcal D\), let
\(\Pi_{\mathcal D}^{\rm out}\) be the orthogonal projection onto the closed
subspace spanned by \(\mathcal D\).  The transition probability into
\(\mathcal D\) is
\[
  P_{\Psi\to\mathcal D}
  =
  \left\|
    \Pi_{\mathcal D}^{\rm out}S\Psi_{\rm in}
  \right\|^2 .
\]
This is the primary definition.  Sharp-momentum formulas are distributional
kernels obtained from this expression after smearing with wave packets and
then taking a narrow-packet limit.

For positive-energy momenta in the mostly-plus metric,
\[
  p_i=(E_i,\vec p_i),
  \qquad
  q_a=(E_a,\vec q_a),
\]
with
\[
  p_i^2=-m_i^2,\qquad q_a^2=-M_a^2,
\]
set
\[
  P_{\rm in}=\sum_i p_i,
  \qquad
  P_{\rm out}=\sum_a q_a .
\]
The invariant amplitude convention inherited from LSZ is
\[
  \langle q_1,\ldots,q_m|T|p_1,\ldots,p_n\rangle
  =
  (2\pi)^D\delta^D(P_{\rm out}-P_{\rm in})
  \mathcal M(q_1,\ldots,q_m|p_1,\ldots,p_n).
\]
The bra and ket are generalized basis vectors in \(\mathcal H_{\rm as}\).
Equivalently, the same distribution is obtained from the connected part of
the Haag--Ruelle kernel
\[
  {}_{\rm out}\langle q_1,\ldots,q_m|p_1,\ldots,p_n\rangle_{\rm in}
\]
after subtracting the identity contribution.  The connected contribution to
\(S\) is therefore
\[
  \ii(2\pi)^D\delta^D(P_{\rm out}-P_{\rm in})\mathcal M .
\]

\section{Wave-Packet Origin of Cross Sections}

Cross sections are transition probabilities per incoming transverse flux.
Let \(\psi_1,\psi_2\in\Hilb_1\) be normalized incoming one-particle packets
with narrow momentum supports and distinct group velocities.  In the
center-of-mass frame let \(\vec p_*\) be the central momentum of the first
packet, and let
\[
  \mathbb R^{D-2}_\perp
  =
  \{b\in\mathbb R^{D-1}: b\cdot \vec p_*=0\}
\]
be the transverse impact-parameter plane.  If \(\tau_b\) denotes the spatial
translation acting on the one-particle Hilbert space, define
\[
  \Psi_b=\psi_1\odot \tau_b\psi_2 .
\]
Here \(\odot\) denotes the two-particle product in the asymptotic Fock space,
with the species labels and Bose or Fermi statistics appropriate to the
sector under consideration.  For an outgoing channel \(\mathcal D\) whose
projection removes the identity no-scattering component, the wave-packet
cross section is
\[
  \sigma_{\psi_1,\psi_2}(\mathcal D)
  =
  \int_{\mathbb R^{D-2}_\perp}\dd^{D-2}b\,
  \left\|
    \Pi_{\mathcal D}^{\rm out}S\Psi_b
  \right\|^2 .
\]
The dimension is length to the power \(D-2\), the dimension of transverse
area in \(D\) spacetime dimensions.

The plane-wave formula is the narrow-packet density of this Hilbert-space
quantity.  The Fourier transform in \(b\) enforces transverse momentum
conservation.  The long-time Haag--Ruelle localization enforces energy and
longitudinal momentum conservation in the scattering limit.  Together these
produce the single overall distribution
\[
  (2\pi)^D\delta^D(P_{\rm out}-P_{\rm in})
\]
appearing in the invariant amplitude convention.  The incoming current
through the transverse impact-parameter plane is the invariant relative speed
multiplied by the one-particle normalizations.  Thus the denominator in the
sharp-momentum two-particle formula is the invariant flux
\[
  \mathcal F(p_1,p_2)=4E_1E_2v_{\rm rel}.
\]
The next section expresses the same factor in covariant form.

\section{Invariant Phase Space and Flux}

For one on-shell final particle of mass \(M\), the Lorentz-invariant
positive-energy measure is
\[
  \dd\mu_M(q)
  =
  \frac{\dd^Dq}{(2\pi)^{D-1}}\,
  \theta(q^0)\delta(q^2+M^2)
  =
  \frac{\dd^{D-1}\vec q}{(2\pi)^{D-1}2E_{\vec q}},
\]
where
\[
  E_{\vec q}=\sqrt{\vec q^{\,2}+M^2}.
\]
For \(m\) final particles with total incoming momentum \(P\), define the
\(m\)-body phase-space measure
\[
  \dd\Phi_m(P;q_1,\ldots,q_m)
  =
  (2\pi)^D\delta^D\!\left(P-\sum_{a=1}^m q_a\right)
  \prod_{a=1}^m\dd\mu_{M_a}(q_a).
\]
If some final particles are identical, the phase-space integral over labeled
momenta is divided by the factorial of the number of identical particles of
each species.

For a two-particle initial state, the invariant flux factor is
\[
  \mathcal F(p_1,p_2)
  =
  4\sqrt{(p_1\cdot p_2)^2-m_1^2m_2^2}.
\]
Equivalently,
\[
  \mathcal F=4E_1E_2v_{\rm rel},
  \qquad
  v_{\rm rel}
  =
  \frac{\sqrt{(p_1\cdot p_2)^2-m_1^2m_2^2}}{E_1E_2}.
\]
In the center-of-mass frame,
\[
  p_1=(E_1,\vec p_*),
  \qquad
  p_2=(E_2,-\vec p_*),
\]
so
\[
  \mathcal F=4(E_1+E_2)|\vec p_*|.
\]

The differential cross section for
\[
  p_1+p_2\longrightarrow q_1+\cdots+q_m
\]
is
\[
  \dd\sigma
  =
  \frac{1}{\mathcal F(p_1,p_2)}
  \left|\mathcal M(q_1,\ldots,q_m|p_1,p_2)\right|^2
  \dd\Phi_m(P_{\rm in};q_1,\ldots,q_m).
\]
All factors in this expression have now been fixed: the asymptotic state
normalization, the amplitude convention, the final-state measure, and the
incoming flux.

\begin{center}
\begin{tikzpicture}[scale=0.86,>=Stealth]
  \draw[blue!75!black,very thick,->] (-4.4,-0.9)--(-1.35,-0.22);
  \draw[green!50!black,very thick,->] (-4.4,0.9)--(-1.35,0.22);
  \draw[purple!80!black,thick,fill=purple!7]
    (-1.15,0) ellipse (0.35 and 1.05);
  \draw[red!75!black,very thick,->] (-0.72,0.28)--(2.05,1.1);
  \draw[red!75!black,very thick,->] (-0.72,0)--(2.15,0.0);
  \draw[red!75!black,very thick,->] (-0.72,-0.28)--(2.0,-1.1);
  \node[blue!75!black] at (-3.45,-1.2) {$p_1,\ \vec v_1$};
  \node[green!50!black] at (-3.45,1.2) {$p_2,\ \vec v_2$};
  \node[purple!80!black,align=center] at (-1.15,-1.55)
    {effective\\area \(\sigma\)};
  \node[red!75!black] at (2.55,1.1) {$q_1$};
  \node[red!75!black] at (2.55,0.0) {$q_2$};
  \node[red!75!black] at (2.55,-1.1) {$q_m$};
  \node[align=center] at (4.1,0)
    {final\\channel};
  \node[align=center] at (-4.15,0)
    {\(\mathcal F=4E_1E_2v_{\rm rel}\)};
\end{tikzpicture}
\end{center}

\section{Two-Body Final States in Four Dimensions}

Set \(D=4\).  For a two-body final state, define
\[
  s=-(p_1+p_2)^2=P_{\rm in}^0{}^2-\vec P_{\rm in}^{\,2}.
\]
In the center-of-mass frame,
\[
  P_{\rm in}=(\sqrt s,\vec0).
\]
The Kallen polynomial is
\[
  \lambda_K(s,a,b)=s^2+a^2+b^2-2sa-2sb-2ab .
\]
For incoming masses \(m_1,m_2\) and final masses \(M_1,M_2\),
\[
  |\vec p_*|
  =
  \frac{\sqrt{\lambda_K(s,m_1^2,m_2^2)}}{2\sqrt s},
  \qquad
  |\vec q_*|
  =
  \frac{\sqrt{\lambda_K(s,M_1^2,M_2^2)}}{2\sqrt s}.
\]
Performing the two-body phase-space integral over the energy-conserving
delta function gives
\[
  \dd\Phi_2
  =
  \frac{1}{16\pi^2}
  \frac{|\vec q_*|}{\sqrt s}\,\dd\Omega ,
\]
where \(\dd\Omega\) is the solid-angle measure of one final momentum in the
center-of-mass frame.  The flux is
\[
  \mathcal F=4\sqrt s\,|\vec p_*|.
\]
Therefore
\[
  \frac{\dd\sigma}{\dd\Omega}
  =
  \frac{1}{64\pi^2s}
  \frac{|\vec q_*|}{|\vec p_*|}
  |\mathcal M(s,\Omega)|^2 .
\]
For identical final particles, the right hand side is divided by \(2\) when
the integral is taken over the full labeled solid angle.

For elastic scattering of identical scalar particles of mass \(m\),
\[
  |\vec q_*|=|\vec p_*|,
  \qquad
  \sqrt s=E_{\rm cm}.
\]
If the tree-level \(\phi^4\) invariant amplitude is
\[
  \mathcal M=-g+O(g^2),
\]
then
\[
  \sigma_{2\to2}
  =
  \frac{g^2}{32\pi s}+O(g^3)
\]
for identical final particles.

\section{Unitarity and the Optical Theorem}

The operator identity
\[
  S^\dagger S=\mathbf 1,\qquad S=\mathbf 1+\ii T,
\]
is equivalent to
\[
  -\ii(T-T^\dagger)=T^\dagger T .
\]
Let \(i\) and \(f\) denote incoming and outgoing asymptotic momentum labels,
including all discrete quantum numbers.  Insert a complete set of outgoing
asymptotic states between \(T^\dagger\) and \(T\).  After removing the common
overall momentum-conservation delta function, the sharp-kernel form is
\[
  -\ii\left[
  \mathcal M(f|i)-\mathcal M(i|f)^*
  \right]
  =
  \sum_X
  \int\dd\Phi_X\,
  \mathcal M(X|f)^*\mathcal M(X|i),
\]
where \(X\) runs over all allowed multiparticle channels and
\(\dd\Phi_X\) includes the corresponding phase-space measure and symmetry
factors.

For \(f=i\), this becomes the optical theorem
\[
  2\,\operatorname{Im}\mathcal M(i|i)
  =
  \sum_X
  \int\dd\Phi_X\,|\mathcal M(X|i)|^2 .
\]
For two incoming particles,
\[
  \sigma_{\rm tot}
  =
  \frac{2\,\operatorname{Im}\mathcal M(i|i)}
  {\mathcal F(p_1,p_2)} .
\]
The amplitude on the left is the forward amplitude: the outgoing labels are
the same as the incoming labels.

\section{Partial Waves in Four Dimensions}

For scalar \(2\to2\) scattering in the center-of-mass frame, rotational
invariance makes the amplitude a function of \(s\) and
\[
  x=\cos\theta=\widehat p_*\cdot\widehat q_* .
\]
For distinguishable particles use the partial-wave expansion
\[
  \mathcal M(s,x)
  =
  16\pi
  \sum_{\ell=0}^{\infty}
  (2\ell+1)a_\ell(s)P_\ell(x),
\]
where \(P_\ell\) is the Legendre polynomial.  The inverse formula is
\[
  a_\ell(s)
  =
  \frac{1}{32\pi}
  \int_{-1}^1\dd x\,
  P_\ell(x)\mathcal M(s,x).
\]
For equal incoming masses in the elastic region, define
\[
  \rho(s)=\frac{2|\vec p_*|}{\sqrt s}.
\]
The partial-wave \(S\)-matrix eigenvalue is
\[
  S_\ell(s)=1+2\ii\rho(s)a_\ell(s).
\]
Elastic unitarity is
\[
  |S_\ell(s)|=1.
\]
Equivalently, there is a real phase shift \(\delta_\ell(s)\) such that
\[
  S_\ell(s)=\ee^{2\ii\delta_\ell(s)}
\]
and
\[
  a_\ell(s)
  =
  \frac{\ee^{2\ii\delta_\ell(s)}-1}{2\ii\rho(s)}
  =
  \frac{\ee^{\ii\delta_\ell(s)}\sin\delta_\ell(s)}{\rho(s)} .
\]
The same condition can be written
\[
  \operatorname{Im}a_\ell(s)=\rho(s)|a_\ell(s)|^2 .
\]
With inelastic channels open,
\[
  |S_\ell(s)|\le1,
\]
with equality precisely in a purely elastic partial wave.

For distinguishable scalar particles, the elastic total cross section is
\[
  \sigma_{\rm el}
  =
  \frac{\pi}{|\vec p_*|^2}
  \sum_{\ell=0}^\infty
  (2\ell+1)|S_\ell(s)-1|^2
  =
  \frac{4\pi}{|\vec p_*|^2}
  \sum_{\ell=0}^\infty
  (2\ell+1)\sin^2\delta_\ell(s).
\]

For identical scalar bosons, the physical two-particle amplitude on unordered
states is Bose symmetric.  With the convention
\[
  \mathcal M_{\rm id}(s,x)
  =
  32\pi
  \sum_{\substack{\ell=0\\ \ell\ {\rm even}}}^{\infty}
  (2\ell+1)a_\ell(s)P_\ell(x),
\]
and the full-solid-angle final-state factor \(1/2\), the elastic cross
section is
\[
  \sigma_{\rm el}^{\rm id}
  =
  \frac{2\pi}{|\vec p_*|^2}
  \sum_{\substack{\ell=0\\ \ell\ {\rm even}}}^{\infty}
  (2\ell+1)|S_\ell(s)-1|^2 .
\]
This is the convention used later for identical-particle high-energy bounds.

\section{Partial-Wave Unitarity Geometry}

Define
\[
  b_\ell(s)=\rho(s)a_\ell(s),
  \qquad
  S_\ell(s)=1+2\ii b_\ell(s).
\]
In an elastic partial wave, \(|S_\ell|=1\).  Writing
\[
  b_\ell=x_\ell+\ii y_\ell,
\]
this condition becomes
\[
  x_\ell^2+\left(y_\ell-\frac12\right)^2=\frac14 .
\]
Thus the elastic amplitude lies on the circle of radius \(1/2\) centered at
\(\ii/2\) in the complex \(b_\ell\)-plane.  Equivalently,
\[
  y_\ell=|b_\ell|^2 .
\]
The phase-shift form
\[
  b_\ell=\ee^{\ii\delta_\ell}\sin\delta_\ell
\]
parametrizes this circle.

When inelastic channels are open, the elastic partial wave is a matrix
element of a unitary matrix on a larger channel space.  It is convenient to
write
\[
  S_\ell=\eta_\ell\ee^{2\ii\delta_\ell},
  \qquad
  0\le\eta_\ell\le1 .
\]
Then \(b_\ell=(S_\ell-1)/(2\ii)\) lies inside the same circle:
\[
  x_\ell^2+\left(y_\ell-\frac12\right)^2\le\frac14,
  \qquad
  y_\ell\ge |b_\ell|^2 .
\]
The difference
\[
  y_\ell-|b_\ell|^2
  =
  \frac{1-\eta_\ell^2}{4}
\]
is the contribution of inelastic final states in that angular-momentum
channel.  This is the partial-wave form of the optical theorem.

\section{Thresholds, Poles, and Resonances}

The same variables organize the analytic structure of scattering.  A
threshold is the value of \(s\) at which a new on-shell intermediate channel
becomes kinematically available.  Across such a threshold, unitarity gives
the discontinuity of amplitudes in terms of phase-space integrals over lower
amplitudes.

A stable bound state in a two-particle channel appears as a pole of the
partial-wave amplitude on the physical sheet below threshold.  A resonance is
represented by a pole reached after analytic continuation through a unitarity
cut.  The precise sheet and residue statements require the analytic
continuation machinery developed later, but the organizing data are already
present here: the \(S\)-operator, phase space, unitarity, and partial waves.
